% =============================================================================
%  CHAPTER 3 — DESIGN / ARCHITECTURE
%  RUBRIC: Execution & Technical Quality (50%). BUDGET: ~6-9 pages.
%
%  >>> Sections 3.1-3.5 are an AI FIRST DRAFT from your PROJECT_CONTEXT.md:
%      edit into your voice, verify details, then I polish.
%  >>> 3.1 principle justifications and 3.6 alternatives are left as [WRITE] —
%      that reasoning is yours and is what the viva probes.
% =============================================================================
\chapter{Design}
\label{ch:design}

\section{Design Goals and Principles}
\label{sec:design-principles}

The plugin is shaped by five principles, each a direct response to a failure mode
of text-generation refactoring identified in Chapter~\ref{ch:background}.

\paragraph{Correctness by construction.} Every write is performed by IntelliJ's
own refactoring engine rather than by emitting modified source text. The result
is therefore correct with respect to the IDE's semantic model --- references,
imports, comments, and string occurrences updated consistently --- independently
of the language model's text-generation quality. This is the project's central
design commitment, and the property that distinguishes it from text-patch agents.

\paragraph{Decoupling reasoning from execution.} The language model decides
\emph{what} to refactor; the IDE decides and performs \emph{how}. The agent never
writes code: it selects an operation and its arguments, and the engine carries it
out. This narrows the agent's responsibility to the task it is good at (planning)
and removes the task it is unreliable at (producing correct edits).

\paragraph{A complete, multi-operation surface.} Rather than a single hardwired
operation, the plugin exposes the full range an agent needs to do realistic
maintenance --- rename, move, change-signature, safe-delete, inline, extract, and
symbol creation --- so that multi-step tasks can be completed end to end without
falling back to text editing for the operations the tool does not cover.

\paragraph{Model-agnosticism.} The interface is a language-neutral JSON tool
schema served over HTTP, and re-exported over the Model Context Protocol, so that
any function-calling model can drive it. The contribution is the structured
editing interface, not a dependency on any one provider.

\paragraph{Read before write.} Inspection tools let the agent examine source and
assess the blast radius of a change before committing to it, supporting genuine
read--plan--act loops rather than blind edits.

Of these five, the first two are the load-bearing commitments. Correctness by
construction and the decoupling of reasoning from execution are what let the approach
hold up \emph{at scale}: because the engine --- not the model --- propagates each
change to every dependent reference, correctness does not rest on the agent
re-deriving the full blast radius of an edit across a large codebase, and the
guarantee that the change is behaviour-preserving holds regardless of how capable the
driving model is or whether it has a build-feedback loop to catch its own mistakes.
The agent is left only with the part it is genuinely reliable at --- choosing which
operation to apply --- which is precisely why relying on the engine's AST-aware view
of the project is sounder than trusting the model's text edits to be complete.

\section{System Architecture}
\label{sec:design-architecture}

Figure~\ref{fig:arch} shows the path a single operation takes. An external agent
--- or the in-IDE chat panel --- issues a tool call as an HTTP \texttt{POST} to a
server bound to \texttt{127.0.0.1}; the server dispatches the call to a
project-scoped service method; that method resolves the target symbol and invokes
the corresponding IntelliJ processor, which mutates the PSI and persists the
result to disk. Because the same backend serves both the external HTTP surface
and the internal chat panel, the two entry points share one implementation and
one correctness guarantee.

\begin{figure}[h]
\centering
\begin{tikzpicture}[
  font=\small,
  box/.style={draw, rounded corners, align=center, inner sep=5pt, minimum width=78mm},
  eng/.style={box, fill=black!7},
  >={Stealth[]},
  node distance=5mm,
]
  \node[box] (agent) {External coding agent (Claude Code or any\\ function-calling LLM), or the in-IDE chat panel};
  \node[box, below=of agent] (http) {\texttt{POST /tools} on \texttt{127.0.0.1}\\ (also re-exported over the Model Context Protocol)};
  \node[box, below=of http] (server) {\texttt{AgentToolServer} $\rightarrow$ \texttt{ToolsHandler.dispatch}};
  \node[box, below=of server] (svc) {\texttt{RefactorService}: resolve qualified name\\ to a \texttt{PsiElement}};
  \node[eng, below=of svc] (proc) {IntelliJ refactoring engine\\ (\texttt{HeadlessRenameProcessor}, \texttt{ChangeSignatureProcessor}, \ldots)};
  \node[eng, below=of proc] (psi) {PSI semantic model: every reference, import\\ and occurrence updated atomically};
  \node[box, below=of psi] (files) {Source files on disk};
  \foreach \a/\b in {agent/http,http/server,server/svc,svc/proc,proc/psi,psi/files}
    \draw[->] (\a) -- (\b);
  \node[right=3mm of proc, font=\footnotesize\itshape, text width=23mm, align=left] {correct by construction};
\end{tikzpicture}
\caption{Request path from an agent's tool call to a correct-by-construction edit.}
\label{fig:arch}
\end{figure}

As a concrete trace, a rename begins as a call naming the target by qualified name
and a new identifier; the service resolves the name to a \texttt{PsiElement},
hands it to \texttt{RenameProcessor}, and the processor updates the declaration
and every resolved reference across the project as a single atomic refactoring.
In code, the request lands in \texttt{ToolsHandler.handle}, which parses the JSON body
and routes it through \texttt{dispatch} to the project-scoped \texttt{RefactorService}
(resolved via \texttt{project.service<RefactorService>()}). For the rename above,
\texttt{RefactorService.renameByQualifiedName} resolves the name to a
\texttt{PsiElement} and runs a \texttt{HeadlessRenameProcessor} --- a
\texttt{RenameProcessor} subclass that suppresses the otherwise-blocking preview and
conflict dialogs (\Cref{ch:impl}) --- whose \texttt{run()} performs the project-wide
update inside a single write action before the change is flushed to disk. The hosting
\texttt{AgentToolServer} registers \texttt{ToolsHandler} at \texttt{/tools} alongside
\texttt{SchemaHandler} (\texttt{/tools/schema}) and \texttt{StatusHandler}
(\texttt{/status}), all bound to \texttt{127.0.0.1}.

\section{The Tool API Surface}
\label{sec:design-tools}

The API exposes twenty-four operations, grouped by role
(Table~\ref{tab:tools}). All are published as a single JSON schema at
\texttt{GET /tools/schema}, which the agent loads to discover what it can do.

\begin{table}[h]
\centering
\caption{The tool surface, grouped by role.}
\label{tab:tools}
\begin{tabular}{@{}lll@{}}
\toprule
Group & Tools & Correctness basis \\
\midrule
Lookup     & find-symbol-by-name, list-symbols, find-symbol & PSI symbol resolution \\
Refactor   & rename, safe-delete, move-class, change-signature & IntelliJ processors \\
Hierarchy / Inline & pull-up-member, push-down-member, make-static, & IntelliJ processors \\
           & migrate-type, inline-field, inline-method & \\
Create (Java)   & add-field, add-method, add-inner-class, create-file & \texttt{PsiElementFactory} \\
Create (Kotlin) & add-property, add-function, create-file & \texttt{KtPsiFactory} \\
Extract    & extract-method, extract-variable & headless \texttt{MethodExtractor} \\
Inspect    & read-file, find-usages & PSI reference index \\
\bottomrule
\end{tabular}
\end{table}

The Hierarchy/Inline group (pull-up and push-down member, make-method-static,
type-migration, and field/method inlining) is implemented and exposed through the same
schema but is not exercised by the benchmark suites of \Cref{ch:eval}, which concentrate
on the rename, move, change-signature, safe-delete, and usage-discovery core; benchmarking
these higher-level operations is noted as future work (\Cref{sec:concl-future}).

The granularity is chosen to match the unit at which IntelliJ actually guarantees
correctness: one tool per \emph{refactoring operation}. A single coarse
\texttt{refactor} tool would push the hard part --- deciding \emph{which} semantic
operation a request implies --- back onto the model, reintroducing the ambiguity the
design exists to remove. A surface of dozens of fine-grained, PSI-node-level edits
would, in the limit, recreate text editing one node at a time and forfeit the
by-construction guarantee, since correctness in IntelliJ is a property of whole
operations (\texttt{RenameProcessor}, \texttt{MoveClassesOrPackagesProcessor},
\texttt{ChangeSignatureProcessor}), not of individual node mutations. The
operation-level surface also matches what the benchmark tasks actually demand: each
task corresponds to one or two named operations, so the agent completes it by issuing
those operations directly rather than by planning a sequence of low-level edits.

\section{Symbol Resolution by Qualified Name}
\label{sec:design-symbols}

Prior agents address code by file position or byte offset, which forces the model
to guess locations and re-guess after each edit. This design instead addresses
symbols by \emph{qualified name} --- for example \texttt{com.example.Foo\#bar} ---
so the agent operates in the semantic domain of names rather than positions.
Overloaded methods are disambiguated by typed signature: a call naming
\texttt{Owner\#getPet(String,boolean)} resolves to exactly one of three
\texttt{getPet} overloads, and the rename touches only that one and its callers.
This is precisely what the \texttt{pc-rename-method-002} task isolates
(\Cref{sec:eval-rootcause}): renaming only the two-argument
\texttt{Owner\#getPet(String,boolean)} overload while the other two \texttt{getPet}
overloads must survive. The three overloads are lexically identical, so a tool that
matches on the \emph{name} alone cannot tell them apart --- the scripted baseline
renames all three and is caught by the surviving-symbol check. Resolving the target by
typed signature instead lets the engine select exactly the intended overload and
rename it together with precisely its callers, by construction; a text-level approach
can reach the same outcome only by reconstructing, from the surrounding code, the
type information that the signature already carries.

\section{Model-Agnostic Access via MCP}
\label{sec:design-mcp}

To make the model-agnostic claim concrete, a bridge re-exports the HTTP tool API
as a Model Context Protocol server over stdio. It fetches \texttt{/tools/schema}
live, so any tool added to the plugin appears automatically to connected clients,
and any MCP-capable agent (for example Claude Desktop or the Gemini CLI) can drive
the same structured backend with no change to the plugin. This realises the core
architectural separation of the project: the language model is a replaceable
planning layer, and the IntelliJ PSI engine is the execution layer.

\section{Design Alternatives}
\label{sec:design-alternatives}
The design followed fairly directly from the gap analysis of
Chapter~\ref{ch:background} rather than from a wide search over competing
architectures. Several alternative designs nonetheless exist at the key decision
points, and it is worth stating why the chosen approach is preferable to each.

\paragraph{Generate a text patch, then validate it.} A natural alternative keeps the
model generating modified source text and adds a compiler or linter pass afterwards to
reject bad patches. This still places unsafe text generation on the write path: a
compile check catches edits that fail to build, but not edits that compile yet are
semantically wrong --- the over-eager overload rename of \Cref{sec:eval-rootcause}
compiles and would pass such a gate. It also presupposes a build-feedback loop, which
the structured engine does not need, since its edits are correct by construction.

\paragraph{Keep the model inside the IDE, with no external interface.} EM-Assist and
similar tools embed the model in the IDE and expose no callable API. That forecloses
the central contribution --- a model-agnostic surface that \emph{any} external
function-calling agent (Claude Code, the Gemini CLI, or a benchmark harness) can drive
--- and ties the system to a single in-IDE assistant. Exposing the engine over HTTP and
MCP is what makes the execution layer reusable and independently testable.

\paragraph{Address code by position rather than by name.} Prior agents locate targets
by file offset, which the model must guess and then re-guess after each edit shifts the
offsets, and which cannot distinguish overloaded methods. Addressing symbols by
qualified, signature-typed name (\Cref{sec:design-symbols}) keeps the agent in a stable
semantic namespace and lets the engine resolve exactly one overload --- the property
the \texttt{pc-rename-method-002} task exercises (\Cref{sec:eval-rootcause}).
